\documentclass{article}

\begin{document}
	
	DECISION TREES — SUMMARY (UNIT 1.4)
	
	\bigskip
	
	\textbf{1. What is a Decision Tree}
	
	A decision tree is a supervised learning model mainly used for classification.
	It represents a decision process as a tree structure composed of internal nodes,
	branches, and leaves. Internal nodes test attributes, branches correspond to
	attribute values, and leaves assign class labels. Classification is obtained
	by following a path from the root to a leaf.
	
	\bigskip
	
	\textbf{2. Logical Interpretation}
	
	Decision trees can be interpreted as rule-based models. Each root-to-leaf path
	corresponds to a conjunction of conditions, while the full tree represents a
	disjunction of conjunctions. This logical form makes decision trees highly
	interpretable and suitable for explanation.
	
	\bigskip
	
	\textbf{3. Space Partitioning View}
	
	From a geometric perspective, decision trees recursively partition the feature
	space into disjoint regions. Each leaf corresponds to one region. The splits are
	usually axis-aligned, meaning that each decision tests a single attribute at a
	time. Deeper trees create more and smaller regions.
	
	\bigskip
	
	\textbf{4. Tree Construction}
	
	Decision trees are built using a top-down greedy strategy. Starting from the
	root, the algorithm selects the best attribute to split the data, partitions
	the dataset accordingly, and repeats the process recursively. The construction
	stops when the data are pure or no attributes remain.
	
	\bigskip
	
	\textbf{5. ID3 Algorithm}
	
	ID3 is a classical algorithm for building decision trees. It constructs the tree
	top-down and selects attributes using information gain. ID3 assumes discrete
	attributes, does not handle continuous features natively, and does not perform
	pruning. As a result, it may produce deep trees that overfit the training data.
	
	\bigskip
	
	\textbf{6. Entropy}
	
	Entropy is a measure of uncertainty or impurity in a dataset. It is low when most
	instances belong to the same class and high when classes are evenly distributed.
	Entropy is minimal for pure nodes and maximal when class proportions are equal.
	
	\bigskip
	
	\textbf{7. Information Gain}
	
	Information gain measures the reduction in entropy achieved by splitting the
	data using a particular attribute. Attributes that produce purer subsets have
	higher information gain. However, attributes with many values may be favored,
	which can bias the splitting process.
	
	\bigskip
	
	\textbf{8. Overfitting in Decision Trees}
	
	Decision trees are prone to overfitting, especially when they grow very deep or
	when the data contain noise. Overfitting leads to high variance and poor
	generalization to unseen data.
	
	\bigskip
	
	\textbf{9. Pruning}
	
	Pruning is used to reduce overfitting by simplifying the tree structure.
	Pre-pruning stops the growth of the tree early, while post-pruning removes
	branches after the full tree has been built. Pruning improves generalization,
	reduces variance, and often increases interpretability.
	
	\bigskip
	
	\textbf{10. Strengths and Limitations}
	
	\textbf{Strengths}
	\begin{itemize}
		\item High interpretability
		\item Natural handling of categorical data
		\item No need for feature scaling
	\end{itemize}
	
	\textbf{Limitations}
	\begin{itemize}
		\item Greedy and locally optimal construction
		\item High variance
		\item Axis-aligned splits only
		\item Sensitivity to noise
	\end{itemize}
	
\end{document}
